% TOP_PROBLEM: {"id": 30007049, "title": "Zero Lelong numbers for minimal metrics on nef adjoints", "category_id": 5, "category": {"id": 5, "name": "algebraic_geometry", "display_name": "Algebraic Geometry", "description": "Geometric objects defined by polynomial equations.", "slug": "algebraic-geometry", "order_index": 5, "created_at": "2026-07-31T15:26:25.671Z"}, "status": "open"}
% FIRST_POSTED: 2026-10-06
% ATTEMPT_STATUS: solved
% OPENAI_CLAIM: {"schema_version":1,"source_repo":"https://github.com/openai/math","source_commit":"adc7f1241b42e322a6451854ab7e4b4c146bf78a","release_date":"2026-10-06","match":"full","resolution":"proved","independently_reviewed":false,"second_pass":"agreed","families":[{"family":"034","title":"Log abundance for compact Kähler spaces under logarithmic Iitaka subadditivity","subject":"Algebraic and complex geometry","manuscripts":[{"title":"Minimal metrics and interior injectivity for nef adjoints","pdf_path":"preprints/Minimal-metrics-and-interior-injectivity-for-nef-adjoints-September-27-2026/paper.pdf","readme_path":"preprints/Minimal-metrics-and-interior-injectivity-for-nef-adjoints-September-27-2026/README.md","date":"2026-09-27","theorem_ref":"thm:adjoint-metric (sections/introduction.tex:29)","statement_sources":[{"path":"preprints/Minimal-metrics-and-interior-injectivity-for-nef-adjoints-September-27-2026/build/sections/introduction.tex","line":29,"label":"thm:adjoint-metric"}]}],"lean":null,"reasoning_trace":null}]}
\documentclass[11pt]{article}
\usepackage[margin=25mm]{geometry}
\usepackage{fontspec}
\setmainfont{Latin Modern Roman}
\usepackage{amsmath,amssymb,mathtools}
\usepackage{unicode-math}
\setmathfont{Latin Modern Math}
\usepackage{enumitem}
\usepackage{xurl,hyperref,bookmark}
\hypersetup{colorlinks=true,urlcolor=blue}
\setcounter{secnumdepth}{0}
\setlength{\parindent}{0pt}
\setlength{\parskip}{5pt}
\setlength{\emergencystretch}{3em}
\begin{document}
\section*{Zero Lelong numbers for minimal metrics on nef adjoints}

\subsection{OpenAI's claim}
Minimal semipositive metrics of nef klt adjoints are claimed to have no positive Lelong numbers, and zero-Lelong endpoint metrics are claimed to give the stated interior-boundary $H^1$ injection. The cited manuscript argues for that result; its relation to this problem is specified below.

\subsection{Mathematical statement of the claimed result}
\paragraph{Mathematical claim 1.} For every normal connected projective complex klt pair $(H,\Theta)$ with effective rational $\Theta$ and nef rational Cartier $D_H=K_H+\Theta$, and every projective log resolution $\pi:V\to H$, OpenAI claims that $N=\pi^*D_H$ admits a semipositive singular Hermitian metric of minimal singularities and
\[
\nu(h,x)=0\quad\text{for every }x\in V
\]
for every such metric $h$. Here $\nu$ is the Lelong number; minimal singularities means its weights are no more singular, up to bounded additive error, than the weights of any other semipositive metric on $N$. The components $K_H$ and $\Theta$ need not be separately rational Cartier.

\subsection{Outline of the argument}
Manuscript 1 (family 034, 2026-09-27): An extremal normalized-section argument contradicts a positive Lelong number by improving the minimal weight locally. For injectivity, harmonic forms and local primitives are compared with both endpoint weights; weak convergence in a fixed Hilbert space retains the ordinary coherent-cohomology class.

\subsection{Scope relative to this problem}
Manuscript 1 (family 034, 2026-09-27): The problem selects the mathematical statement summarized here, including its stated hypotheses and parameter range. The cited theorem gives a claimed proof of that assertion in full.

\subsection{Formal verification}
Family 034 has no Lean scope page in this snapshot; no formalization of these manuscript results is listed.

This site has not executed the Lean files or checked their dependencies. These descriptions report the scope of the linked declarations, rather than a new verification of the proofs.

\subsection{Status caveat}
OpenAI attributes these results to an internal model. Its release notice cautions that results lacking full formalization can contain mistakes. This site has not independently assessed the proofs; a solved label records the scope of the reported claim.

\subsection{Sources}
\begin{itemize}
\item OpenAI, \emph{Minimal metrics and interior injectivity for nef adjoints}, 2026-09-27; citation key \texttt{OAI:Minimal-metrics-and-interior-injectivity-for-nef-adjoints-September-27-2026}. \href{https://github.com/openai/math/blob/adc7f1241b42e322a6451854ab7e4b4c146bf78a/preprints/Minimal-metrics-and-interior-injectivity-for-nef-adjoints-September-27-2026/paper.pdf}{PDF (GitHub)}; \href{https://github.com/openai/math/raw/adc7f1241b42e322a6451854ab7e4b4c146bf78a/preprints/Minimal-metrics-and-interior-injectivity-for-nef-adjoints-September-27-2026/paper.pdf}{PDF (direct)}; \href{https://github.com/openai/math/blob/adc7f1241b42e322a6451854ab7e4b4c146bf78a/preprints/Minimal-metrics-and-interior-injectivity-for-nef-adjoints-September-27-2026/README.md}{Manuscript page and citation}; \href{https://github.com/openai/math/blob/main/preprints/Minimal-metrics-and-interior-injectivity-for-nef-adjoints-September-27-2026/paper.pdf}{Latest version}; Statement source: \href{https://github.com/openai/math/blob/adc7f1241b42e322a6451854ab7e4b4c146bf78a/preprints/Minimal-metrics-and-interior-injectivity-for-nef-adjoints-September-27-2026/build/sections/introduction.tex\#L29}{thm:adjoint-metric}.
\item \href{https://github.com/openai/math/blob/adc7f1241b42e322a6451854ab7e4b4c146bf78a/CONTENTS.md}{OpenAI catalogue entry}.
\end{itemize}
\end{document}
